\documentclass{subfiles}
\begin{document}
    Let \(x\) be an integer and let \(B(x)\) be it's binary expansion,
        i.e. it's binary representaion. Then, the gamma encoding of \(x\)
        consist of two parts: the unary representation of \(\abs{B(x)}\) and 
        the binary representation of \(x\) with the most significant digit removed.
        That is 
        \[
            \gamma(x) = U(\abs{B(x)}) \tilde{B}(x)
        \]
        where \(\tilde{B}(x)\) is \(B(x)\) with the most significant digit removed.

    \begin{example*}
        Consider the integer 29. Its binary representation \(B(29) = 11101\).
            The unary encoding of \(\abs{B(29)}\) is \(U(\abs{B(29)}) = U(5) = 00001\).
            Therefore, the gamma encoding of \(29\) is
            \[
                \gamma(29) = \underbrace{00001}_{U(\abs{B(29)})}
                    \underbrace{1101}_{\tilde{B}(29)}
            \]
    \end{example*}

    Decoding is trivial: count the number of zeros up to the first 1,
        say this are \(k\), fetch the following \(k + 1\) bits and interpret the 
        sequence as the integer \(x\). 

    It can be proved that gamma encoding is optimal when the distibution is 
        \[
            p(x) = \frac{1}{x^{2}}.
        \]
        As easily seen, the inefficiency is given by the unary encoding of \(\abs{B(x)}\).
\end{document}
